\subsection{The Weierstrass Extreme Value Theorem}

\begin{definition}
  Let $D \subseteq \mathbb{R}$ and $f : D \rightarrow \mathbb{R}$.
  \begin{itemize}
    \item We say $c \in D$ is a {\bf minimum point} of $f$ in $D$ if 
      \[f(c) \leq f(x) \quad \forall x \in D.\]
  \item We say $c \in D$ is a {\bf maximum point} of $f$ in $D$ if
      \[f(c) \geq f(x) \quad \forall x \in D .\]
  \end{itemize}
  We say $f$ {\bf attains} its maximum or minimum in $D$ if a maximum or minimum point exists in $D$.
\end{definition}

\begin{theorem}
    Let $a, b \in \mathbb{R}$ with $a < b$. Suppose $f : [a, b] \rightarrow \mathbb{R}$ is continuous. Then
    \begin{itemize}
      \item there exists finite $L > 0$ such that
        \[|f(x) | \leq L \quad \forall x \in [a, b].\]
      \item $f$ attains its maximum and minimum in $[a, b]$.
    \end{itemize}
\end{theorem}
\begin{proof} $ $
  \begin{itemize}
    \item By contradiction, we have for any $n$, there exists an $x \in [a,b]$ such that
      \[f(x) > n.\]
      Define $(x_n)_{n \in \mathbb{N}}$ in this way. Then this sequence is clearly bounded so has a convergent subsequence $x_{n_k} \rightarrow c$. As $f$ is continuous
      \[f(c) = \lim_{n \rightarrow \infty} f(x_{n_k}) = \infty\]
      which is a contradiction as $f(c)$ is clearly bounded.
    \item Let $M = \sup \{f(x) | x \in [a, b]\}$ which is bounded. Then for all $n$ there exists a $y$ such that
      \[M - \frac{1}{n} \leq y \leq M\]
      and define the sequence $y_n$ in this way, so $y_n \rightarrow M$. Since $y$ is bounded it has a convergent subsequence $y_{n_k} \rightarrow c \in [a,b]$ and also by continuity
      \[f(c) = \lim_{n \rightarrow \infty} f(y_{n_k}) = M.\]
      hence $f$ obtains a maximum. The same argument can be used for the supremum.
  \end{itemize}    
\end{proof}

\subsection{The Natural Logarithm}

Before we can define the inverse of the exponential function, we must show that it is a bijection.
\begin{prop} The exponential function is bijective.
\end{prop}
\begin{proof} We first prove injectivity and then surjectivity.
  \begin{itemize}
    \item  Let $x, y \in \mathbb{R}$ with $x > y$. Then
    \begin{align*}
      \frac{\exp(x)}{\exp(y)} &= \exp(x - y) \\
                              &= 1 + \sum_{n=1}^{\infty} \frac{(x - y)^n}{n!} > 1
    .\end{align*}
    Therefore, as the sum is positive, $\exp(x) > \exp(y)$. This implies that $\exp : \mathbb{R} \rightarrow (0, \infty)$ is injective.
  \item Fix $y > 0$. Then
    \[\exp(y) = 1 + y + \sum_{n=2}^{\infty} \frac{y^n}{n!} > y.\]
    This implies, for $y > 0$, $\exp(\frac{1}{y}) > \frac{1}{y} \implies \frac{1}{\exp(y)} = \exp(- \frac{1}{y}) < y$.

    Hence, $\exp(- \frac{1}{y}) < y < \exp(y)$, so there is a $c$ such that $\exp(c) = y$.
  \end{itemize}
\end{proof}

\begin{definition}

  The function $\log : (0, \infty) \rightarrow \mathbb{R}$ is the inverse of $\exp : \mathbb{R} \rightarrow (0, \infty)$.
    
\end{definition}

\begin{prop}
    For $x, y \in (0, \infty)$
    \begin{itemize}
      \item $\log(xy) = \log x + \log y$
      \item $\log( \frac{x}{y}) = \log x - \log y$
    \end{itemize}
\end{prop}
\begin{proof} $ $
  \begin{itemize}
    \item For the first proof
      \begin{align*}
        \exp(\log x + \log y) &= \exp(\log x)\exp(\log y) \\
                               &= xy \\
                               &= \exp(\log(x y))
    .\end{align*}
    so by injectivity $\log xy = \log x + \log y$
  \item Because $\exp(0) = 1$ this means $\log 1 = 0$.
    \begin{align*}
      \log( \frac{y}{y}) &= \log y + \log( \frac{1}{y}) \\
                      0  &= \log y + \log( \frac{1}{y})\\
                      \log( \frac{1}{y}) &= -\log(y)
    \end{align*}
    and the result follows.
  \end{itemize}
\end{proof}

\begin{definition}
  For $a > 0$ and $x \in \mathbb{R}$
  \[a^x = \exp(x \log a).\]
\end{definition}

\section{Derivatives}

\begin{definition}
Suppose $S \subseteq \mathbb{R}$. A point $c \in S$ is called an {\bf interior point} of $S$ if there exists $\delta > 0$ such that
\[(c - \delta, c + \delta) \subseteq S.\]
\end{definition}

\begin{definition}
  Let $D \subseteq \mathbb{R}$ and $f: D \rightarrow \mathbb{R}$. Suppose $c \in D$ is an interior point of $D$. We say $f$ is {\bf differentiable} at $c$ if
  \[\lim_{h \rightarrow 0} \frac{f(c + h) - f(c)}{h} \quad \text{exists}.\]

  The value $f'(c) := \lim_{h \rightarrow 0} \frac{f(c + h) - f(c)}{h}$ is called {\bf the derivative} of $f$ at $c$.

  If $f$ is differentiable at every point $c \in D$, we say $f$ is {\bf differentiable} on $D$.
\end{definition}

\begin{prop}
    Let $D \subseteq \mathbb{R}$ and let $c \in D$ be an interir point of $D$. If $f:D \rightarrow \mathbb{R}$ is differentiable at $c$, then it is continuous at $c$.
\end{prop}
\begin{proof}
  Suppose $\lim_{h \rightarrow 0} \frac{f(c + h) - f(c)}{h}$ exists. Using a substition of $x = c + h$, this can be rewritten as
  \begin{align*}
    \lim_{x \rightarrow c} \frac{f(x) - f(c)}{x - c} &= f'(c) \\
    \lim_{x \rightarrow c} f(x) - f(c) &= \lim_{x \rightarrow c} (x - c) f'(c) \\
    \lim_{x \rightarrow c} f(x) - f(c) &= 0
  \end{align*}
  so $f$ is continuous at $c$.
\end{proof}

\subsection{The Algebra of Derivatives}

\begin{theorem}
    Let $D \subseteq \mathbb{R}$ and let $c \in D$ be an interior point. Suppose that $f, g : D \rightarrow \mathbb{R}$ are differentiable at $c$. Then
    \begin{itemize}
      \item $f + g$ is differentiable at $c$ with
        \[(f + g)'(c) = f'(c) + g'(c).\]
      \item $f - g$ is differentiable at $c$ with 
        \[(f - g)'(c) = f'(c) - g'(c).\]
      \item $\alpha f$ is differentiable at $c$ with $(\alpha f)'(c) = \alpha f'(c)$
      \item (Product rule) $fg$ is differentiable at $c$ with
        \[(fg)'(c) = f'(c)g(c) - f(c)g'(c).\]
      \item (Quotient rule). If $g(c) \not = 0$ then $\frac{f}{g}$ is differentiable at $c$ with
        \[\left(\frac{f}{g}\right)'(c) = \frac{f'(c)g(c) - f(c)g'(c)}{(g(c))^2}.\]
    \end{itemize}
\end{theorem}
\begin{proof}
    The first three points can be proven relatively easily using the algebra of function limits.
    \begin{itemize}
      \item First, it is easy to show that
        \[(fg)(c + h) - (fg)(c) = \left(f(c + h) - f(c)\right)g(c + h) + f(c)\left(g(c + h) - g(c)\right).\]
        Then
        \begin{align*}
          (fg)'(c) &= \lim_{h \rightarrow 0}\frac{\left(f(c + h) - f(c)\right)g(c + h) + f(c)\left(g(c + h) - g(c)\right)}{h} \\
                   &= f'(c) g(c) + f(c) g'(c)
        .\end{align*}
      \item We only need to prove that $\left( \frac{1}{g}\right)'(c) = \frac{-g'(c)}{(g(c))^2}$, and can then use the product rule. Using the continuity of $g$ with $\epsilon = \frac{|g(c)|}{2}$, we can find a $\delta$ such that for all $h \in (c - \delta, c + \delta)$, $g(c + h) \not = 0$, so it is well defined. Then
        \begin{align*}
          \lim_{h \rightarrow 0} \frac{ ( \frac{1}{g})(c + h) - (\frac{1}{g})(c)}{h} &= \lim_{h \rightarrow 0} \frac{ \frac{1}{g(c +h)} - \frac{1}{g(c)}}{h} \\
           &= - \lim_{h \rightarrow 0} \frac{g(c + h) - g(c)}{h} \cdot \frac{1}{g(c + h) g(c)} \\
           &= - \frac{g'(c)}{g(c)^2}
        .\end{align*}
    \end{itemize}
    
\end{proof}
